\documentclass[10pt,a4paper]{amsart}
\usepackage[margin=19mm]{geometry}
\usepackage{amsmath,amssymb,mathtools}
\usepackage[scr=rsfs,cal=euler]{mathalfa}
\usepackage{xfrac}
\usepackage[unicode,hidelinks,bookmarks=false]{hyperref}
\newcommand{\ie}{i.e.\ }
\newcommand{\cf}{cf.\ }
\newcommand{\resp}{respectively\ }
\newcommand{\Subsection}{\subsection}
\providecommand{\nobreakdash}{\nobreak\hbox{-}}
\setlength{\emergencystretch}{2em}
\multlinegap0pt
\usepackage{amssymb}
\newtheorem{lemma}{Lemma}
\title{Difference of weighted composition operators between some spaces of analytic functions}
\author{Jiaoye Du, Cezhong Tong, Zicong Yang}
\date{Source: arXiv:2507.19735v2 (CC BY 4.0)}
\begin{document}
\maketitle

\noindent\emph{Source and licence.} Difference of weighted composition operators between some spaces of analytic functions. Version arXiv:2507.19735v2 (CC BY 4.0), distributed by arXiv under the Creative Commons Attribution 4.0 International licence, http://creativecommons.org/licenses/by/4.0/, as recorded in the arXiv licence metadata for this exact version and shown on the abstract page https://arxiv.org/abs/2507.19735v2. Full licence notice: \texttt{licenses/arxiv-2507.19735-CC-BY-4.0.txt}.\par\smallskip

\noindent\textit{Editorial scope.} Counted unit: the article's lemma lemma2.3 (source0.tex lines 352--358). The article's complete original proof of it is delivered with it (source0.tex lines 552--628, one \texttt{proof} environment of 77 lines, which the article places away from the statement). No mathematics is written, repaired or re-derived: the statement and the proof are the article's own lines, in the article's own notation, and no retained line is substituted or re-flowed.  Retained uncounted as the source-local context the proof needs: lines 274--276; lines 278--280; lines 282--284; lines 424--436; lines 505--511; lines 513--515.  Nothing was omitted from any retained span, so the package carries no omission record.  No retained line contains a citation, so the wrapper carries no bibliography and no bibliography entry.  No retained line contains a figure environment, an image-inclusion command or drawing code, so no image is delivered and the package declares no asset dependency.  Editorial formatting only, in addition: an AMS \texttt{amsart} preamble that restates the article's own declarations and macros used by the retained lines, namely the environments \texttt{lemma} on separate counters, together with the standard packages those lines need.  The retained environments print in this document as Lemma 1 in order of appearance, whereas the article prints the same items with the numbering of its own full document.  The complete native source of the article as distributed in the arXiv e-print for this exact version is bundled unchanged as \texttt{sources/182238/source0.tex}.
\begin{equation}\label{equa2.2}
A_{\alpha}(D(z,r))\simeq (1-|z|^2)^{2+\alpha},
\end{equation}

\begin{equation}\label{equa2.3}
1-|z|^2\simeq 1-|w|^2\simeq |1-\overline{z}w|
\end{equation}

\begin{equation}\label{equa2.4}
|1-\overline{a}z|\simeq |1-\overline{a}w|
\end{equation}

\begin{lemma}\label{lemma2.7}
Let $0<p,q<\infty$ and $\mu$ be a positive Borel measure on $\mathbb{D}$.
\begin{itemize}
\item[(i)] If $0<p\leq q<\infty$, then $\mu$ is a $(p,q)$-Bergman Carleson measure if and only if 
$$\sup_{z\in\mathbb{D}}M_{r,\frac{q}{p}}(z)<\infty$$
for some (or any) $r>0$, and $\mu$ is a vanishing $(p,q)$-Bergman Carleson measure if and only if
$$\lim_{|z|\to 1^-}M_{r,\frac{q}{p}}(z)=0$$
for some (or any) $r>0$.
\item[(ii)] If $0<q<p<\infty$, then $\mu$ is a $(p,q)$-Bergman Carleson measure if and only if it is a vanishing $(p,q)$-Bergman Carleson measure if and only if
$$M_r(\mu)\in L^{\frac{p}{p-q}}(\mathbb{D},{\rm d}A_{\alpha})$$
for some (or any) $r>0$.
\end{itemize}
\end{lemma}

\begin{equation}\label{equa3.4}
\begin{split}
&|u(z)-v(z)|\left|\frac{1-\overline{a}\varphi(z)}{1-\overline{a}\psi(z)}\right|^N\\
&\quad \leq \left|u(z)-v(z)\left(\frac{1-\overline{a}\varphi(z)}{1-\overline{a}\psi(z)}\right)^N\right|+|u(z)|\left|1-\left(\frac{1-\overline{a}\varphi(z)}{1-\overline{a}\psi(z)}\right)^N\right|\\
&\quad \lesssim \left|u(z)-v(z)\left(\frac{1-\overline{a}\varphi(z)}{1-\overline{a}\psi(z)}\right)^N\right|+|u(z)|\left|\frac{\varphi(z)-\psi(z)}{1-\overline{a}\psi(z)}\right|
\end{split}
\end{equation}

\begin{equation}\label{equa3.5}
\left|\frac{1-\overline{a}\varphi(z)}{1-\overline{a}\psi(z)}\right|\gtrsim \frac{(1-|\varphi(z)|^2)(1-|\psi(z)|^2)}{|1-\overline{\varphi(z)}\psi(z)|^2}=1-\rho^2(z)
\end{equation}

\begin{lemma}\label{lemma2.3}
Let $\{a_j\}$ be an $r$-lattice in the Bergman metric. Suppose $\alpha>-1$, $0<p<\infty$ and $N,i$ be non-negative integers with $N>\max\{1,\frac{1}{p}\}+\frac{1+\alpha}{p}$. For any $\mathbf{c}=\{c_j\}\in l^p$, let
\begin{equation*}
H_{\mathbf{c}, N,p}^{[i]}(z)=\sum_{j=1}^{\infty}c_j\frac{(1-|a_j|^2)^{N+i-\frac{2+\alpha}{p}}z^i}{(1-\overline{a}_jz)^{N+i}},\quad z\in\mathbb{D}.
\end{equation*}
Then $H_{\mathbf{c},N,p}^{[i]}\in A_{\alpha}^p(\mathbb{D})$ and $\|H_{\mathbf{c}, N,p}^{[i]}\|_{A_{\alpha}^p}\lesssim \|\{c_j\}\|_{l^p}$.
\end{lemma}

\begin{proof}[{\bf New proof of $(i)\Rightarrow (iii)$ in Theorem 1.2}]
Assume $0<q<p<\infty$ and $C_{u,\varphi}-C_{v,\psi}: A_{\alpha}^p(\mathbb{D})\to A_{\alpha}^q(\mathbb{D})$ is bounded. Fix $r>0$. Let $\{a_j\}$ be an $r$-lattice in the Bergman metric and $\mathbf{c}=\{c_j\}\in l^p$. For $N>\max\{1,\frac{1}{p}\}+\frac{1+\alpha}{p}$, we use the test function $H_{\mathbf{c},N,p}^{[i]}$ defined in Lemma \ref{lemma2.3} to obtain
\begin{equation*}
\|(C_{u,\varphi}-C_{v,\psi})H_{\mathbf{c},N,p}^{[i]}\|_{A_{\alpha}^q}\lesssim \|\{c_j\}\|_{l^p}, \quad i=0,1.
\end{equation*}
For $i=0$, we obtain
\begin{equation}\label{equa3.6}
\begin{split}
&\int_{\mathbb{D}}\left|\sum_{j=1}^{\infty}c_j\left(u(z)\frac{(1-|a_j|^2)^{N-\frac{2+\alpha}{p}}}{(1-\overline{a}_j\varphi(z))^N}-v(z)\frac{(1-|a_j|^2)^{N-\frac{2+\alpha}{p}}}{(1-\overline{a}_j\psi(z))^N}\right)\right|^q{\rm d}A_{\alpha}(z)\\
&\quad \lesssim \|\{c_j\}\|_{l^p}^q.
\end{split}
\end{equation}
In \eqref{equa3.6}, we replace $c_j$ by $r_j(t)c_j$, so that the right hand side does not change. Then intergrating both sides with respect to $t$ from 0 to 1, and by Khinchine's inequaltiy and Fubini's Theorem, we obtain
\begin{equation*}
\begin{split}
&\int_{\mathbb{D}}\left(\sum_{j=1}^{\infty}|c_j|^2\left|u(z)\frac{(1-|a_j|^2)^{N-\frac{2+\alpha}{p}}}{(1-\overline{a}_j\varphi(z))^N}-v(z)\frac{(1-|a_j|^2)^{N-\frac{2+\alpha}{p}}}{(1-\overline{a}_j\psi(z))^N}\right|^2\right)^{\frac{q}{2}}{\rm d}A_{\alpha}(z)\\
&\simeq\int_{\mathbb{D}}\int_0^1\left|\sum_{j=1}^{\infty}c_jr_j(t)\left(u(z)\frac{(1-|a_j|^2)^{N-\frac{2+\alpha}{p}}}{(1-\overline{a}_j\varphi(z))^N}-v(z)\frac{(1-|a_j|^2)^{N-\frac{2+\alpha}{p}}}{(1-\overline{a}_j\psi(z))^N}\right)\right|^q{\rm d}t{\rm d}A_{\alpha}(z)\\
&=\int_{0}^1\int_{\mathbb{D}}\left|\sum_{j=1}^{\infty}c_jr_j(t)\left(u(z)\frac{(1-|a_j|^2)^{N-\frac{2+\alpha}{p}}}{(1-\overline{a}_j\varphi(z))^N}-v(z)\frac{(1-|a_j|^2)^{N-\frac{2+\alpha}{p}}}{(1-\overline{a}_j\psi(z))^N}\right)\right|^q{\rm d}A_{\alpha}(z){\rm d}t\\
&\lesssim \|\{c_j\}\|_{l^p}^q.
\end{split}
\end{equation*}
By \eqref{equa2.3}, $\chi_{D(a_j,2r)}(\varphi(z))\lesssim \frac{1-|a_j|^2}{|1-\overline{a}_j\varphi(z)|}$. It follows that
\begin{equation*}
\begin{split}
\int_{\mathbb{D}}\left(\sum_{j=1}^{\infty}\frac{|c_j|^2\chi_{D(a_j,2r)}(\varphi(z))\left|u(z)-v(z)\left(\frac{1-\overline{a}_j\varphi(z)}{1-\overline{a}_j\psi(z)}\right)^N\right|^2}{(1-|a_j|^2)^{2(2+\alpha)/p}}\right)^{q/2}{\rm d}A_{\alpha(z)} \lesssim \|\{c_j\}\|_{l^p}^q.
\end{split}
\end{equation*}
Recall that there is a positive integer $N_0$ such that each point $z\in\mathbb{D}$ belongs to at most $N_0$ of the disks $D(a_j,2r)$. Applying Minkowski's inequality if $\frac{2}{q}\leq 1$ and H\"older's inequality if $\frac{2}{q}>1$, we obtain
\begin{equation*}
\begin{split}
&\sum_{j=1}^{\infty}|c_j|^q\frac{\int_{\varphi^{-1}(D(a_j,2r))}\left|u(z)-v(z)\left(\frac{1-\overline{a}_j\varphi(z)}{1-\overline{a}_j\psi(z)}\right)^N\right|^q{\rm d}A_{\alpha}(z)}{(1-|a_j|^2)^{q(2+\alpha)/p}}\\
&\quad\leq C\int_{\mathbb{D}}\left(\sum_{j=1}^{\infty}\frac{|c_j|^2\chi_{D(a_j,2r)}(\varphi(z))\left|u(z)-v(z)\left(\frac{1-\overline{a}\varphi(z)}{1-\overline{a}\psi(z)}\right)^N\right|^2}{(1-|a_j|^2)^{2(2+\alpha)/p}}\right)^{q/2}{\rm d}A_{\alpha(z)}\\
&\quad \lesssim \|\{c_j\}\|_{l^p}^q,
\end{split}
\end{equation*}
where $C=\max\{N_0^{1-\frac{q}{2}},1\}$. This implies that the sequence 
\[\left\{\frac{\int_{\varphi^{-1}(D(a_j,2r))}\left|u(z)-v(z)\left(\frac{1-\overline{a}_j\varphi(z)}{1-\overline{a}_j\psi(z)}\right)^N\right|^q{\rm d}A_{\alpha}(z)}{(1-|a_j|^2)^{q(2+\alpha)/p}}\right\}_{j=1}^{\infty}\]
belongs to the dual of $l^{\frac{p}{q}}$ or, equivalently,
\begin{equation}\label{equa3.7}
\sum_{j=1}^{\infty}\left(\frac{\int_{\varphi^{-1}(D(a_j,2r))}\left|u(z)-v(z)\left(\frac{1-\overline{a}_j\varphi(z)}{1-\overline{a}_j\psi(z)}\right)^N\right|^q{\rm d}A_{\alpha}(z)}{(1-|a_j|^2)^{q(2+\alpha)/p}}\right)^{\frac{p}{p-q}}<\infty.
\end{equation}
Similarly, taking $i=1$, we obtain
\begin{equation}\label{equa3.8}
\sum_{j=1}^{\infty}\left(\frac{\int_{\varphi^{-1}(D(a_j,2r))}\left|u(z)\varphi(z)-v(z)\psi(z)\left(\frac{1-\overline{a}_j\varphi(z)}{1-\overline{a}_j\psi(z)}\right)^{N+1}\right|^qdA_{\alpha}(z)}{(1-|a_j|^2)^{q(2+\alpha)/p}}\right)^{\frac{p}{p-q}}<\infty.
\end{equation}
By \eqref{equa2.3}, $\left|\frac{1-\overline{a}_j\varphi(z)}{1-\overline{a}_j\psi(z)}\right|\lesssim 1$ when $\varphi(z)\in D(a,r)$ and constant suppressed is independent of $a_j$, multiplying the integrand in \eqref{equa3.7} by $|\psi(z)|^q\left|\frac{1-\overline{a_j}\varphi(z)}{1-\overline{a}_j\psi(z)}\right|^q$, adding it to \eqref{equa3.8} and by the triangle inequality and \eqref{equa2.4}, we get
\begin{equation}\label{equa3.9}
\sum_{j=1}^{\infty}\left(\frac{\int_{\varphi^{-1}(D(a_j,2r))}|u(z)|^q\rho(z)^q{\rm d}A_{\alpha}(z)}{(1-|a_j|^2)^{q(2+\alpha)/p}}\right)^{\frac{p}{p-q}}<\infty.
\end{equation}
Therefore, by \eqref{equa2.2} and \eqref{equa2.3}, we obtain
\begin{equation}\label{equa3.10}
\begin{split}
&\int_{\mathbb{D}}\left(\frac{\int_{\varphi^{-1}(D(w,r))}|u(z)|^q\rho(z)^q{\rm d}A_{\alpha}(z)}{(1-|w|^2)^{2+\alpha}}\right)^{\frac{p}{p-q}}{\rm d}A_{\alpha}(w)\\
&\quad \leq \sum_{j=1}^{\infty}\int_{D(a_j,r)}\left(\frac{\int_{\varphi^{-1}(D(w,r))}|u(z)|^q\rho(z)^q{\rm d}A_{\alpha}(z)}{(1-|w|^2)^{2+\alpha}}\right)^{\frac{p}{p-q}}{\rm d}A_{\alpha}(w)\\
&\quad \lesssim \sum_{j=1}^{\infty}\left(\frac{\int_{\varphi^{-1}(D(a_j,2r))}|u(z)|^q\rho(z)^q{\rm d}A_{\alpha}(z)}{(1-|a_j|^2)^{2+\alpha}}\right)^{\frac{p}{p-q}}A_{\alpha}(D(a_j,2r))\\
&\quad \lesssim \sum_{j=1}^{\infty}\left(\frac{\int_{\varphi^{-1}(D(a_j,2r))}|u(z)|^q\rho(z)^q{\rm d}A_{\alpha}(z)}{(1-|a_j|^2)^{q(2+\alpha)/p}}\right)^{\frac{p}{p-q}}<\infty.
\end{split}
\end{equation}
This implies that $\omega_{\varphi,u}^q$ is a $(p,q)$-Bergman Carleson measure by Lemma \ref{lemma2.7}. Similarly, $\omega_{\psi,v}^q$ is also a $(p,q)$-Bergman Carleson measure.

On the other hand, when $\varphi(z)\in D(a_j,2r)$, by \eqref{equa3.4}, \eqref{equa3.5} and \eqref{equa2.4}, we have
\begin{equation*}
\begin{split}
|u(z)-v(z)|(1-\rho(z))^N&\lesssim |u(z)-v(z)|\left|\frac{1-\overline{a}_j\varphi(z)}{1-\overline{a}_j\psi(z)}\right|^N\\
&\lesssim \left|u(z)-v(z)\left(\frac{1-\overline{a}_j\varphi(z)}{1-\overline{a}_j\psi(z)}\right)^N\right|+|u(z)|\rho(z).
\end{split}
\end{equation*}
So we combine \eqref{equa3.7} and \eqref{equa3.9} to obtain
\begin{equation*}
\sum_{j=1}^{\infty}\left(\frac{\int_{\varphi^{-1}(D(a_j,2r))}|u(z)-v(z)|^q(1-\rho(z))^{qN}{\rm d}A_{\alpha}(z)}{(1-|a_j|^2)^{q(2+\alpha)/p}}\right)^{\frac{p}{p-q}}<\infty.
\end{equation*}
Then through a similar argument as \eqref{equa3.10}, we get
\begin{equation*}
\int_{\mathbb{D}}\left(\frac{\int_{\varphi^{-1}(D(w,r))}|u(z)-v(z)|^q(1-\rho(z))^{qN}{\rm d}A_{\alpha}(z)}{(1-|w|^2)^{2+\alpha}}\right)^{\frac{p}{p-q}}{\rm d}A_{\alpha}(w)<\infty.
\end{equation*}
This implies that $\sigma_{\varphi,\beta}^q$ is a $(p,q)$-Bergman Carleson measure for $\beta\geq qN>\frac{q}{p}(\alpha+1+\max\{1,p\})$ by Lemma \ref{lemma2.7}. Similarly, $\sigma_{\psi,\beta}^q$ is also a $(p,q)$-Bergman Carleson measure. The proof is now complete.
\end{proof}

\end{document}
